\band{orangehead}{Practice Midterm · Part 2}
\k{Part 2 รวม 25 คะแนน}: วางแผนเป็น PM; grooming เป็น PO; sprint planning เป็น SM; ไม่ตัดสินแทนทีม
\topic{1) DormFix: WBS 3 + RACI 2 + Gantt 5}
\k{บริบท}: เว็บแจ้งซ่อมหอพักแบบ Waterfall; เราเป็น PM; Aom = BA, Bank = UX/UI, Chai = Dev, Dao = Dev/Tester; ผจก.หอ = sponsor, หัวหน้าช่าง = user
\k{งานและกำหนดการ}: d = วันทำงาน; ทุกวันด้านล่างอยู่ใน Jan 2027; ช่องท้าย = วันทำ / วันเผื่อ
{\smalltable\begin{tabularx}{\linewidth}{|C{2.6mm}|Y|C{8.5mm}|}\hline
\rowcolor{lightrow}ID & งาน · ผู้ทำ · d & ทำ / เผื่อ\\\hline
E & เขียนข้อกำหนดเป็นงานแรก · Aom · 3d & 4-6 / –\\\hline
B & ส่งข้อกำหนดให้ ผจก.หออนุมัติ · PM · 1d & 7 / +8\\\hline
C & ออกแบบจอจากข้อกำหนดอนุมัติ · Bank · 3d & 11-13 / –\\\hline
G & ออกแบบ DB จากข้อกำหนดอนุมัติ · Chai · 2d & 11-12 / –\\\hline
F & ส่งแบบจอ+DB ให้หัวหน้าช่างอนุมัติ · PM · 1d & 14 / +15\\\hline
A & โค้ดฟอร์มแจ้งซ่อมจากแบบอนุมัติ · Chai · 4d & 18-21 / –\\\hline
I & โค้ด dashboard ช่างจากแบบอนุมัติ · Dao · 3d & 18-20 / –\\\hline
D & ทดสอบโปรแกรมสองส่วนร่วมกัน · Dao · 2d & 22,25 / –\\\hline
H & ส่งระบบทดสอบแล้วให้ ผจก.หอตรวจรับ · PM · 1d & 26 / +27\\\hline
PB & Project buffer ท้ายโครงการ · 2d & – / 28-29\\\hline
\end{tabularx}}
\k{ลำดับ}: \b{E → B → (C∥G) → F → (A∥I) → D → H → PB}
{\smalltable\begin{tabularx}{\linewidth}{|C{4mm}|*{5}{>{\centering\arraybackslash}X|}}\hline
\rowcolor{lightrow}Task & PM & Aom & Bank & Chai & Dao\\\hline
E & A & R & C & I & –\\\hline
C & A & C & R & I & –\\\hline
G & A & – & C & R & I\\\hline
D & A & I & – & C & R\\\hline
\end{tabularx}}
\extra RACI จาก notes ของ mock 5.2 \r{ไม่ใช่เฉลยอาจารย์}; “–” = ไม่มีบทบาท
\k{กฎ}: เริ่ม จ. 4 Jan; เต็มเวลา จ.-ศ. ไม่มีวันหยุดอื่น; งานถัดไปเริ่มวันทำงานหลังงานก่อน+buffer จบ
\begin{center}\begin{tikzpicture}[x=.985mm,y=1mm,font=\fontsize{6.2}{6.4}\selectfont]
\draw[line width=.35pt](0,0)rectangle(48,-27.6);
\foreach\i in{0,...,20}{\draw[black!55,line width=.2pt]({6+2.1*\i},0)--({6+2.1*\i},-27.6);}
\foreach\i in{1,...,11}{\draw[black!55,line width=.2pt](0,{-2.3*\i})--(48,{-2.3*\i});}
\foreach\d[count=\i from0]in{4,5,6,7,8,11,12,13,14,15,18,19,20,21,22,25,26,27,28,29}{\node at({7.05+2.1*\i},-1.15){\d};}
\foreach\lab[count=\j from1]in{E,B,C,G,F,A,I,D,H,PB,M}{\node at(3,{-1.15-2.3*\j}){\lab};}
\foreach\row/\start/\end in{1/0/2,2/3/3,3/5/7,4/5/6,5/8/8,6/10/13,7/10/12,8/14/15,9/16/16}{\foreach\i in{\start,...,\end}{\fill[keyblue!25]({6+2.1*\i},{-2.3*\row})rectangle({8.1+2.1*\i},{-2.3*(\row+1)});\node at({7.05+2.1*\i},{-1.15-2.3*\row}){X};}}
\foreach\row/\i in{2/4,5/9,9/17,10/18,10/19}{\node[keyred]at({7.05+2.1*\i},{-1.15-2.3*\row}){–};}
\foreach\i in{4,9,19}{\node[keygreen]at({7.05+2.1*\i},-26.45){*};}
\end{tikzpicture}\end{center}
\k{M}: milestone วันที่ 8/15/29 Jan; X = ทำงาน, – = buffer, * = milestone (0 วัน); Start/Due = วันแรก/ท้ายที่ทำเครื่องหมายรวม buffer
\g{เหตุผล}: B/F/H รอคนนอกที่มีงานอื่น จึงเผื่อหลังงานละ 1 วัน; ทีมภายในใช้ project buffer ร่วม 2 วันท้าย; milestone \r{ไม่เพิ่มวันทำงาน}
\topic{WBS แบบภาพ: กลุ่มงาน ≠ dependency}
\begin{center}\begin{tikzpicture}[x=1mm,y=1mm,font=\fontsize{7.1}{7.4}\selectfont]
\node[act,fill=bluehead!30](root)at(24,0){DormFix};
\foreach\num/\id/\x/\title/\tasks in{1/1.1/7/Req approved/{E\\B},2/1.2/24/Designs approved/{C ∥ G\\F},3/1.3/41/System accepted/{A ∥ I\\D\\H}}{
\node[act,fill=orangehead!22,text width=13mm](m\num)at(\x,-5.5){\id\\\title};
\node[draw,align=center,text width=12.5mm,inner sep=1pt,anchor=north](t\num)at(\x,-9.5){\tasks};
\draw(root.south)--(m\num.north);\draw(m\num.south)--(t\num.north);
}
\end{tikzpicture}\end{center}
\k{ตรวจ WBS}: ใช้ A-I ครั้งเดียว; 1.4 ว่าง; เลขใบ 1.1.1, 1.1.2 ฯลฯ ตามกลุ่ม; ต้นไม้ไม่ใช่ dependency
\topic{2) Backlog grooming (10): Issue \#12}
\b{ในฐานะนักศึกษา ฉันต้องการจองสนามแบดหนึ่งชั่วโมง เพื่อไม่ต้องต่อคิวที่เคาน์เตอร์} (ใช้โครง As a / I need / So that)
\k{รายละเอียด}: มีสนาม/เวลาเปิดแล้ว; นักศึกษาถือการจองได้คนละหนึ่งรายการต่อวัน; Court 2 วันนี้ 17:00 ว่าง
\par\g{Given} Court 2 วันนี้ 17:00 ว่าง\par
\g{When} นักศึกษาจอง Court 2 เวลา 17:00\par
\g{Then} ช่วงนั้นแสดง “Booked”
\extra AC ลบ: มี booking วันนี้อยู่แล้ว → ขอเพิ่ม → ปฏิเสธและไม่สร้างรายการ; ถามเงื่อนไขยกเลิกก่อนขยายกฎ
\k{Metadata}: Assignees –; Labels \b{enhancement}; Milestone –; Estimate –; \r{ยังไม่ planning อย่ากำหนดคน/points/sprint เอง}
\k{Bug variant (Issue \#15)}: DB บันทึกการยกเลิกแล้วแต่สนามยัง Booked คนอื่นจองไม่ได้ → label \b{bug}; Given Court 1 เวลา 18:00 ถูกจอง; When ยกเลิก; Then “Free”; อีก 3 metadata ใช้ “–”
\k{ถามเพิ่ม}: เขตเวลา? ยกเลิกแล้วคืนสิทธิ์วันนั้นไหม? \r{ไม่เติมเงื่อนไขที่ยังไม่ตกลงเป็น AC บังคับ}
\topic{3) Sprint planning (5): Goal + capacity}
\k{ShuttleNow Sprint 6}: goal = นักศึกษาติดตามรถ/รู้เวลาถึง; anchor “ดูรายการเส้นทางรถรับส่ง” = 3 points; Sprint ก่อนเสร็จ 16 คง velocity เดิม; backlog เรียง priority
{\smalltable\begin{tabularx}{\linewidth}{|C{3mm}|Y|C{4mm}|C{5mm}|}\hline
\rowcolor{lightrow}\# & Story / เทียบ anchor & Pt & Pick\\\hline
1 & Live position (+2 ขั้น) & 8 & Yes\\\hline
2 & QR fare (เท่ากัน) & 3 & No\\\hline
3 & Arrival at my stop (+1) & 5 & Yes\\\hline
4 & Arrival alert (เท่ากัน) & 3 & Yes\\\hline
5 & Arrival at all stops (+1) & 5 & No\\\hline
6 & Fix mobile truncation (−1) & 2 & No\\\hline
\end{tabularx}}
\b{เลือก 1/3/4: 8+5+3=16}; 2 ไม่ตรง goal; เพิ่ม 5 รวม 21 เกิน capacity; 6 ไม่มีที่เหลือ. \r{Fibonacci นับขั้น ไม่บวกเลขตรง}
\begin{center}\begin{tikzpicture}[x=1mm,y=1mm,font=\fontsize{7.1}{7.4}\selectfont]
\node[act](a)at(5,0){2};\node[act,fill=bluehead!45](b)at(18,0){3};\node[act](c)at(31,0){5};\node[act](d)at(44,0){8};
\draw[flow](b)--(a);\draw[flow](b)--(c);\draw[flow](c)--(d);
\node[anchor=south,inner sep=.5pt]at([yshift=.6mm]a.north){−1 step};\node[anchor=north,inner sep=.5pt]at([yshift=-.6mm]b.south){anchor};\node[anchor=south,inner sep=.5pt]at([yshift=.6mm]c.north){+1 step};\node[anchor=north,inner sep=.5pt]at([yshift=-.6mm]d.south){+2 steps};
\end{tikzpicture}\end{center}
\topic{PERT (Lecture 6): คำนวณ effort ไม่ใช่วันนัดส่ง}
\k{ข้อมูล O/M/P (h), Lecture 6}: T1 guard 2/3/6; T2 queue 3/5/10; T3 workflow 4/6/11; T4 comments 4/7/13; T5 admin 3/5/12; T6 auth/tests 4/6/14
\g{T1}: E=3.33 h, SD=0.67 h; รวม E=35.67 h, SD=3.23 h เมื่อประมาณการอิสระ → \mbox{35.7 ± 3.2 h} (หนึ่ง SD) \r{ไม่ใช่ช่วงรับประกัน}; วันจบขึ้นกับงานขนาน/คน
\topic{Git: คำสั่ง remote ที่ควรรู้}
\cl{git fetch origin}
รับ commits/อัปเดต origin/*; \r{ไม่เลื่อน local branch, HEAD หรือ working tree} ในการ fetch ปกติ
\cl{git switch main; git merge --ff-only origin/main}
อัปเดต main แบบ FF; ถ้าแยกสายจะปฏิเสธ ไม่สร้าง merge commit
\cl{git push origin sprint\_1\_staging}
เผยแพร่ staging/อัปเดต tracking ref เมื่อสำเร็จ; push feature \r{ไม่เผยแพร่ merge ที่อยู่เฉพาะ staging}
